\section{Source Syntax and Grammar}\label{sec:source-grammar}

Motivo is defined by a lexical scanner and a context-free parser; the authoritative specification is
\texttt{compiler/docs/grammar.bnf}.
The scanner recognizes note literals (\verb|[A-G][#b]?[0-8]|), drum names, keywords for structure
(\texttt{track}, \texttt{pattern}, \texttt{voice}), control flow, and declarable types (\texttt{int},
\texttt{double}, \texttt{bool}, \texttt{note}, \texttt{seq}, \texttt{chord}).
\texttt{rest} and drum literals appear in expressions but are not declaration types.

A program begins with optional \texttt{tempo} and \texttt{signature} declarations, then top-level variables,
patterns, and tracks.
Tracks may contain statements, local patterns, and \texttt{voice} blocks; voices may not nest further tracks.
Declarations use \texttt{<type> IDENT = <expr>;} and pattern parameters use the same type keywords.
\texttt{play} targets notes, rests, chords, sequences, drum hits, or pattern calls, with optional \texttt{:duration}
and \texttt{from} offsets.

\topic{Version 1.0 to 2.0.}
Version~1.0 used untyped \texttt{let} bindings, arity-only pattern overloading, and a second semantic pass that
simulated pattern bodies at call sites.
Version~2.0 replaces \texttt{let} with explicit typed declarations, resolves overloads by full parameter type
signature, and performs semantic analysis in a single pass---because explicit types make parameter contracts
visible in source text and remove the need to infer intent from initializer shape alone.
